\tinysidebar{\debug{exercises/{test/question.tex}}}
\begin{ex}
Using a triple for-loop, initialize all values in the above x to 3.14.

The 2-dimensional array helps us model data that is inherently
2-dimensional. For instance:

\begin{consolethree}
const int FLOOR_SIZE = 5;

const int NUM_ROOMS = 20;

int max_occupancy[FLOOR_SIZE][NUM_ROOMS];

models the maximum occupancy for rooms in a building with 5 floors and
20 rooms per floor. The data is inherently 2-dimensional. For instance
max_occupancy[3][19] is the maximum occupancy of ROOM319 in the
building.

But three dimensional data occurs too. For instance if you want to model
the presence of physical bodies in a 3-dimensional grid of 1-by-1-by-1
cubic feet of space, you can have a 3-dimensional array of booleans:

const int MAX_X = 100;

const int MAX_Y = 100;

const int MAX_Z = 100;

bool has_body[MAX_X][MAX_Y][MAX_Z];
\end{consolethree}

\end{ex}
